\section{Regression}
\subsection{Linear Regression}
In Linear Regression we learn a mapping $\mathbf{f_w}:\mathbb{R}^{m}\mapsto
\mathbb{R}$ between input attributes $x$ and a single output $y$,
drawn from a function class $\mathcal{F}$:
\[\mathcal{F}=\left\{ \mathbf{f_w(X)=w\cdot x|w}={\left[ w_0,w_1,\ldots,w_m \right]}^{T}\in \mathbb{R}^{m+1} \right\}.\] 
$\mathbf{f_w}$ is characterised by a weight vector $\mathbf{w}$ of which we define:
\[\mathbf{f_w(x)=w\cdot x,\quad x}={[1,x_1,\ldots,x_m]}^{T}\]
\subsubsection{Training Data}
Suppose we draw $n$ observations, each consisting $m$ features (each $i$-th
observation encoded $x^{(i)}\in \mathbb{R}^{m}$), as well as the output labels
$y_i\in \mathbb{R}$. We represent our output training data as $y$:
\[\mathbf{y}=\begin{bmatrix} y^{(1)}\\y^{(2)}\\\vdots\\y^{(n)} \end{bmatrix}.\] 
\begin{definition}[Design Matrix]
  To represent the input attribute data more compactly, we define the design matrix $X$:
  \[\mathbf{X}=\begin{bmatrix}
    \mathbf{x}^{(1)T}\\\mathbf{x}^{(2)T}\\\vdots\\\mathbf{x}^{(n)T}
    \end{bmatrix} =\begin{bmatrix}
  1&x_1^{(1)}&\cdots&x_m^{(1)}\\
  1&x_2^{(2)}&\cdots&x_m^{(2)}\\
  \vdots&\vdots&\ddots&\vdots\\
  1&x_1^{(n)}&\cdots&x_m^{(n)}
\end{bmatrix}.\] 
\end{definition}
We therefore seek $\mathbf{w}$ such that $\mathbf{y=Xw}$. We have a solution
only if $rank(\mathbf{X})=rank(\mathbf{X|y})=(m+1)\le n$. In general this is
not true.

\subsubsection{Evaluation}
To make the problem well-defined, we choose an evaluation function
to help us select the optimal weight vector. We consider the empirical squared
error evaluation function:
\[L(\mathcal{E},\mathcal{S},\mathbf{f_w})=\frac{1}{2n}\sum_{i=1}^{n} {(y^{(i)}-\mathbf{w\cdot x^{(i)}})}^{2}.\] 
We seek the $\mathbf{w}$ at which this evaluation function is minimised.
\[\mathbf{w}_{OLS}=\argmin_\mathbf{w}\left(\frac{1}{2}\sum_{i=1}^{n} {(y^{(i)}-\mathbf{w\cdot x^{(i)}})}^{2}\right).\] 
The optimal ${\mathbf{w}}_{OLS}$ is known as the Ordinary Least Squares
estimator and this approach is known as OLS regression.
Assuming the features and the output are related by an additive noise model:
\[y^{(i)}=\mathbf{w\cdot x}^{(i)}+\epsilon^{(i)}.\]
The Gauss Markov Theorem demonstrates that $\mathbf{w}_{OLS}$ is the Best
Linear Unbiased Estimator (BLUE).
Further assuming the noise is normally distributed:
\[\epsilon\sim\mathcal{N}(0,\sigma^{2})\quad \implies\quad
\mathbf{y|x}\sim\mathcal{N}(\mathbf{w\cdot x}, \sigma^{2})\] 
The Maximum Likelihood Estimator $\mathbf{w}_{MLE}$ is:
\begin{align*}
  \mathbf{w}_{MLE}=&\argmax_\mathbf{w}{\left\{\ln\left( \prod_{i=1}^{n} \frac{1}{\sqrt{2\pi\sigma^{2}}}\exp\left(-\frac{{(y^{(i)}-w\cdot x^{(i)})}^{2}}{2\sigma^{2}}\right)  \right)\right\} }\\
  =&\argmax_\mathbf{w}{\sum_{i=1}^{n} \ln\left( \frac{1}{\sqrt{2\pi\sigma^{2}}}\exp\left(-\frac{{(y^{(i)}-w\cdot x^{(i)})}^{2}}{2\sigma^{2}}\right)  \right) }\\
  =&\argmax_\mathbf{w}{\left(-n \ln \sqrt{2\pi\sigma^{2}} - \sum_{i=1}^{n} \left(\frac{{(y^{(i)}-w\cdot x^{(i)})}^{2}}{2\sigma^{2}}\right)  \right) }\\
  =&\argmin_\mathbf{w}{\left( \frac{1}{2}\sum_{i=1}^{n} {\left(y^{(i)}-\mathbf{w\cdot x}^{(i)}\right)}^{2} \right) }\\
  =&\mathbf{w}_{OLS}
\end{align*}

\subsection{Linear Regression --- Analytic Solution}
\subsubsection{Critical Points}
We first find the gradient of the loss function:
\begin{align*}
  L(\mathcal{E},\mathcal{S},f_w)=&\frac{1}{2n}{\left\lVert \mathbf{y-Xw} \right\rVert_2}^2\\
  =&\frac{1}{2n}\mathbf{{(y-Xw)}^{T}(y-Xw)}\\
  \nabla_w L =& \frac{1}{2n}\nabla_w \left(\mathbf{{(y-Xw)}^{T}(y-Xw)}\right)\\
  =& \frac{1}{2n}\nabla_w\left( \mathbf{y^{T}y-y^{T}Xw-w^{T}X^{T}y+w^{T}X^{T}Xw}\right)\\
  =&\frac{1}{2n}\nabla_w\left(\mathbf{w^{T}X^{T}Xw-2y^{T}Xw}\right) \\
  =&\frac{1}{n} \left( \mathbf{X^{T}Xw-X^{T}y} \right)
\end{align*}
and solve for the criticial points of the function:
\begin{align*}
  \nabla_w L =& 0\\
  0=&\frac{1}{n} \left( \mathbf{X^{T}Xw}_{OLS}-\mathbf{X^{T}y} \right)\\
  \mathbf{X^{T}Xw}_{OLS}=&\mathbf{X^{T}y}
\end{align*}
Since $rank(\mathbf{X^{T}X})=rank(\mathbf{X^{T}X|X^{T}y})$, the system is
consistent and we have one or infinitely many solutions.\\
If ${(\mathbf{X^{T}X})}^{-1}$ exists we have the unique solution:
\[\mathbf{w}_{OLS}=\mathbf{(X^{T}X)}^{-1}\mathbf{X^{T}y}.\]
known as the \textbf{Normal Equations} with an analytic solution to the OLS approach.\\
\begin{itemize}
  \item \textbf{Representation} \[\mathcal{F}=\left\{ \mathbf{w\cdot x}|\mathbf{w}\in \mathbb{R}^{m+1} \right\}\]
  \item \textbf{Evaluation} \[\frac{1}{2}\sum_{i=1}^{n} {(y^{(i)}-\mathbf{w\cdot x^{(i)}})}^{2}\]
  \item \textbf{Optimisation} \[\mathbf{w}_{OLS}=\argmin_\mathbf{w}\left\{
      \frac{1}{2}\sum_{i=1}^{n} {(y^{(i)}-\mathbf{w\cdot x^{(i)}})}^{2} \right\}\]
\end{itemize}
If ${(\mathbf{X^{T}X})}$ is non-invertible then $rank(\mathbf{X^{T}X})<(m+1)$,
the problem is undetermined and we have an infinite number of solutions. This
is known as overfitting.
\begin{figure}[h!]
  \centering
  \includegraphics[width=0.4\textwidth]{overfitting}
\end{figure}
\subsubsection{Optimality}
To check for the optimality of the solution, we first obtain the Hessian:
\[\nabla_w^{2}L=\mathbf{X^{T}X}\]
Since the Hessian matrix is positive semidefinite
($\mathbf{a^{T}X^{T}Xa={\left\lVert Xa \right\rVert_2}^{2}}\ge 0\quad \forall
\mathbf{a}\in \mathbb{R}^{m+1}$), L is convex and a critical point is globally
optimal
\subsubsection{Uniqueness}
Uniqueness of the solution requires that Hessian is positive definite
$\iff\mathbf{(X^{T}X)}^{-1}$, such that the function is strictly convex.
Otherwise, the function is not strictly convex (but still bounded since
$\mathbf{X^{T}y}\in range(\mathbf{X^{T}X})$).
\begin{remark}
  Matrix Derivation Results w.r.t. $w$:
  \begin{itemize}
    \item Linear form: \(a^T w\implies\nabla_w (a^T w) = a\) for some constant vector \(a\)
    \item Quadratic form: \(w^T A w \implies\nabla_w (w^T A w) = (A + A^T) w\)for some constant matrix \(A\)
  \end{itemize}
\end{remark}
\subsubsection{Limitations}
The problem does not scale well with the dimensionality of the input attributes
$m$. Matrix inversion requires $O(m^{3})$ operations --- which is expensive for
large $\mathbf{X^{T}X}$. There are also space constraints for very high $m$ to
store $\mathbf{X^{T}X}$.

\subsection{Linear Regression --- Numerical Optimisation}
To overcome these limitations, we consider a first order iterative numerical
optimisation algorithm, \textbf{Gradient Descent}. Intuitively, we take steps proportional
to the negative of the gradient (which is in the direction of steepest descent)
at each step.
\subsubsection{Batch Gradient Descent}
\begin{lstlisting}
def LinearRegression_OLS_GD(X, y, w_init, alpha, max_epoch=10):
    n = len(y)
    w_t = w_init
    L_hist = (1/(2*n)) * (y + np.matmul(X, w_t)).T.dot(y - np.matmul(X, w_t))
    w_hist = w_t.T
    epoch = 0
    while epoch < max_epoch:
        grad = (1/n) * X.T.dot(X.dot(w_t) - y).reshape(np.shape(X)[1], 1)
        w_t = (w_t - alpha * grad)
        loss = (1/(2*n)) * (y - X.dot(w_t)).T.dot(y - X.dot(w_t))[0][0]
        L_hist = np.vstack((L_hist, loss))
        w_hist = np.vstack((w_hist, w_t.T))
        epoch += 1
    return w_hist, L_hist
\end{lstlisting}
$\alpha>0$ is the learning rate which determines the step size (factor of
proportionality to the gradient).For $\alpha$ too small, convergence is slow,
and for $\alpha$ too large, we get divergence.\\
Since the loss function is a convex objective, gradient descent guarantees the
local optimum found is a global optimum. However, Batch Gradient Descent may be
costly, as we need to scan through the entire training set for each step.
\begin{figure}[h!]
  \centering
  \includegraphics[width=0.7\textwidth]{learning-rate}
\end{figure}
\subsubsection{Stochastic Gradient Descent}
\begin{figure}[h!]
  \centering
  \includegraphics[width=0.4\textwidth]{sgd}
\end{figure}
In Stochastic Gradient Descent, we randomly shuffle the dataset for each step,
and estimate the gradient using one observation. Hence, each SGD update is less
optimal than a BGD update, and it will not converge monotonically, but overall
convergence is generally faster.
\subsubsection{Data Scaling}
For multi-dimensional data, even if we make an optimal setting of the learning
rate, \textbf{ill-conditioning} of the loss surface may result in a slow
convergence. This is caused by the gradient of the largest attribute dominating
the overall gradient. We may alleviate the effect through data scaling to scale our attributes. This
ensures that all parameters are updated in similar proportions, leading to
faster convergence.

\subsection{Polynomial Regression}
In Polynomial Regression we learn a mapping $\mathbf{f_W}:\mathbb{R}^{m}\mapsto
\mathbb{R}$ between input attributes $x$ and an output $y$ drawn from a function
class $\mathcal{F}$.$\mathbf{f_w}$ is characterised by a weight vector $\mathbf{w}$:
\[\mathbf{f_w}=w\cdot \hat{z}\]
for some vector $\hat{z}$ to encode the attributes (with higher degree
polynomials and cross-terms).\\
The design matrix $X$ grows with the dimensionality of the input data and the
degree of the polynomial.
\begin{example}[Polynomial of second degree with $2$ attributes]
\[X=\begin{bmatrix} 1&x_1&x_2&x_1^2&x_2^{2}&x_1 x_2 \\
\vdots&\vdots&\vdots&\vdots&\vdots&\vdots\end{bmatrix} \] 
\end{example}
In general, the total number of terms (or multinomial coefficients) is
$\binom{k+p-1}{p-1}$ for $p$ attributes in the $k$-th degree polynomial.

\subsection{Polynomial Regression --- OLS}
We can generate a solution via the normal equations with an analogue to the
normal equations.
\begin{itemize}
  \item \textbf{Representation} \[\mathcal{F}=\left\{ \mathbf{w\cdot x} \right\}\]
  \item \textbf{Evaluation} \[\frac{1}{2}\sum_{i=1}^{n} {(y^{(i)}-\mathbf{w\cdot x^{(i)}})}^{2}\]
  \item \textbf{Optimisation} \[\mathbf{w}_{OLS}=\argmin_\mathbf{w}\left\{
      \frac{1}{2}\sum_{i=1}^{n} {(y^{(i)}-\mathbf{w\cdot x^{(i)}})}^{2} \right\}\]
\end{itemize}

\subsubsection{Overfitting vs. Underfitting}
Underfitting occurs when we restrict ourselves to a representation that is not
rich enough to encompass the relationships. This is characterised by high bias,
and can be remedied by extending our representation to accommodate greater
functional complexity.\\
While increasing functional complexity, we have the scope to overfit the data.
This occurs when the model fits the residual variation (the noise) and contains
more parameters that can be justified by the data. This results in poor
generalisation. It can be characterised by high variance (between empirical and
generalisation loss).
\subsubsection{Learning Curves}
To prevent underfitting and overfitting, we may compare how well the function
performs on test data and training data:
\begin{figure}[h!]
  \centering
  \includegraphics[width=0.4\textwidth]{learning-curve}
\end{figure}
\subsubsection{Bias-Variance Tradeoff}
We will never be able to completely reduce both effects in the absence of
complete population data. As we increase the complexity of the function class
the bias decreases but the variance increases as the hypothesis space grows.

The OLS analytical solution is given by the normal equations
$\mathbf{w}=\mathbf{X^{T}X}^{-1}\mathbf{X^{T}y}$. For $X\in \mathbb{R}^{n\times
k+1}$, $\mathbf{X^{T}X}$ is non-invertible if $k+1>n$ and we do not have a
unique solution.
\begin{remark}
  With a high enough degree polynomial we can always reduce our training error
  to zero.
\end{remark}
This can be remedied through regularisation, to manage model complexity.
Intuitively, this reduces potential for larger magnitude weights related to
higher complexity, and shrinks some of the weights towards zero.

\subsubsection{Regularisation}
In adoption of a Bayesian perspective, we model $\mathbf{w}$ as a random
variable, and the MAP solution to weight optimisation results in a regularised
objective. Depending on the distribution we select for $\mathbf{w}$, we may
derive ridge regression or LASSO.\\
With the PAC approach, we model a worst case probabilistic bound on the
generalised evaluation function $E_\mathcal{D}[L]$ given some function class.
We may also derive ridge regression or LASSO depending on which function class
is selected.

\subsection{Polynomial Regression --- Ridge Regression}
In Ridge Regression we optimise the empirical squared error loss function subject
to a constraint on $\left\lVert \mathbf{w} \right\rVert_2 \le t$.
\begin{itemize}
  \item \textbf{Representation} \[\mathcal{F}=\left\{ \mathbf{w\cdot x}|{\left\lVert \mathbf{w} \right\rVert}_2^2\le t \right\}\]
  \item \textbf{Evaluation} \[\frac{1}{2}\sum_{i=1}^{n} {(y^{(i)}-\mathbf{w\cdot x^{(i)}})}^{2}\]
  \item \textbf{Optimisation} \[\mathbf{w}_{OLS}=\argmin_\mathbf{w}\left\{
  \frac{1}{2}\sum_{i=1}^{n} {(y^{(i)}-\mathbf{w\cdot x^{(i)}})}^{2} +\frac{\lambda}{2}\sum_{i=1}^{m+1} {\left(w^{(i)}\right)}^2\right\}\]
\end{itemize}
Using Lagrange Multipliers, our optimisation problem is now as follows:
\[\nabla_\mathbf{w}\mathcal{L}(\mathbf{w},\lambda)=\mathbf{0}\]
\begin{align*}
  \text{subject to:}\quad&\left\lVert \mathbf{w} \right\rVert_2^2-t\le 0\\
                         &\lambda\ge 0\\
                         &\lambda\left( \left\lVert \mathbf{w} \right\rVert_2^2 -t \right) =0
\end{align*}
where $\mathcal{L}(\mathbf{w},\lambda)=\frac{1}{2}\sum_{i=1}^{n} \left( y^{(i)}-{\mathbf{w\cdot x}^{(i)}}^{2} \right) +\frac{\lambda}{2}\left( \left\lVert \mathbf{w} \right\rVert_2^2-t \right) $\\
If $\lambda=0$, the constraint is redundant, and the solution is the
non-regularised one. If $\lambda>0$, $\left\lVert \mathbf{w} \right\rVert_2
^2=t$ and we have an equality constraint.\\
For any given $t$, we can find the optimal $\lambda^{*}$ by solving for
$\nabla_\lambda \mathcal{L}(\mathbf{w},\lambda^{*})=\mathbf{0}$. Additionally,
$\mathcal{L}$ is convex w.r.t. $\mathbf{w}$ and we achieve optimum at
$\nabla_\mathbf{w} \mathcal{L}\left( \mathbf{w},\lambda \right) =\mathbf{0}$.
The new evaluation function then corresponds to this new objective:
\[\tilde{L}\left( \mathcal{E},\mathcal{S},\mathbf{f_w} \right)=\frac{1}{2}\left\lVert \mathbf{y-Xw} \right\rVert_2 ^{2}+\frac{\lambda}{2}\left\lVert \mathbf{w} \right\rVert_2 ^{2}\]
We can then solve for optimal $\mathbf{w}_{RR}$:
\begin{align*}
  \nabla_\mathbf{w}\tilde{L}=&\mathbf{X^{T}Xw-X^{T}y+\lambda w}\\
  \mathbf{w}_{RR}=&{\left( \mathbf{X^{T}X+\lambda I} \right) }^{-1}\mathbf{X^{T}y}.
\end{align*}
Since the Hessian $\mathcal{H}=\left( \mathbf{X^{T}X+\lambda I} \right) \succeq
0$ and is always invertible, $\tilde{L}$ is strictly convex and we obtain an
unique solution (with gradient descent).
\subsubsection{Complexity Tuning}
The regularisation parameter $\lambda$ and polynomial degree $k$ together are
complexity tuning hyperparameters:
\begin{itemize}
  \item High $\lambda$, low $k\implies$ low functional complexity $\implies$ underfitting
  \item Low $\lambda$, high $k \implies$ high functional complexity $\implies$ overfitting
\end{itemize}
However, in ridge regression, with a large number of weights, most of the
regularised weights can shrink to almost zero but will rarely reach zero. This
makes model interpretation difficult, especially with large number of
parameters.

\subsection{Polynomial Regression --- LASSO}
In LASSO, we optimise for the empirical squared error loss function, now
subject to a contraint on $\left\lVert \mathbf{w} \right\rVert_1 \le t$. This
has the effect of shrinking some weights more aggressively to zero.
\begin{itemize}
  \item \textbf{Representation} \[\mathcal{F}=\left\{ \mathbf{w\cdot x}|\left\lVert \mathbf{w} \right\rVert_1 \le t\right\}\]
  \item \textbf{Evaluation} \[\frac{1}{2}\sum_{i=1}^{n} {(y^{(i)}-\mathbf{w\cdot x^{(i)}})}^{2}\]
  \item \textbf{Optimisation} \[\mathbf{w}_{OLS}=\argmin_\mathbf{w}\left\{
    \frac{1}{2}\sum_{i=1}^{n} {(y^{(i)}-\mathbf{w\cdot x^{(i)}})}^{2}
+\frac{\lambda}{2}\sum_{i=1}^{m+1} {\left(w^{(i)}\right)}\right\}\]
\end{itemize}
However the evaluation function
\[\tilde{L}\left( \mathcal{E},\mathcal{S},\mathbf{f_w} \right)=\frac{1}{2}\sum_{i=1}^{n} {\left( \mathbf{y-Xw} \right)}^{T}\left( \mathbf{y-Xw} \right) +\frac{\lambda}{2}\left\lVert \mathbf{w} \right\rVert_1\] 
is non-differentiable and we cannot look for stationary points in order to
solve. We consider the subgradients:
\[\begin{cases}
  \frac{\partial L}{\partial w_j} +\frac{\lambda}{2} & \text{for } w_j>0\\
  \frac{\partial L}{\partial w_j} -\frac{\lambda}{2}& \text{for }w_j<0
\end{cases}\]
At $w_j=0$ we are at an optimal point if
\begin{align*}
  \frac{\partial L}{\partial w_j} +\frac{\lambda}{2}>0\quad & \text{and;}\quad \frac{\partial L}{\partial w_j} -\frac{\lambda}{2}<0\\
  \implies \frac{\lambda}{2}< &\frac{\partial L}{\partial w_j} <\frac{\lambda}{2}
\end{align*}
\subsubsection{Limitations}
\begin{itemize}
  \item When dealing with high dimensional data with few examples ($n<k+1$), LASSO
  selects at most $n$ variables before it saturates.
  \item With a highly correlated set of variables, LASSO tends to select one from the group and 
  ignore the others.
  \item LASSO optimisation is not necessarily strictly convex, and in general will not
  yield an unique solution.
\end{itemize}

\subsection{Polynomial Regression --- Elastic Net}
In Elastic Net Regularisation, Ridge Regression and LASSO is combined to remedy
various shortcomings.
\begin{itemize}
  \item \textbf{Representation} \[\mathcal{F}=\left\{ \mathbf{w\cdot x}|{\left\lVert \mathbf{w} \right\rVert}_2^2\le t_1 \wedge\left\lVert \mathbf{w} \right\rVert_1 \le t_2\right\}\]
  \item \textbf{Evaluation} \[\frac{1}{2}\sum_{i=1}^{n} {(y^{(i)}-\mathbf{w\cdot x^{(i)}})}^{2}\]
  \item \textbf{Optimisation} \[\mathbf{w}_{OLS}=\argmin_\mathbf{w}\left\{
  \frac{1}{2}\sum_{i=1}^{n} {(y^{(i)}-\mathbf{w\cdot x^{(i)}})}^{2} +\frac{\lambda}{2}\left( \alpha \sum_{i=1}^{m+1} {\left(w^{(i)}\right)}^2+(1-\alpha)\sum_{i=1}^{m+1} w^{(i)} \right) \right\}\]
  for some $0\le \alpha\le 1,\lambda\ge 0$
\end{itemize}
\subsubsection{Beneficial Properties}
\begin{itemize}
  \item Objective is strictly convex, and optimal Elastic Net weights are unique.
  \item The 1-norm term promotes sparsity.
  \item The 2-norm term promotes a grouping effect and removes the limitation
    on the number of selected variables.
\end{itemize}
